\documentclass[12pt,letterpaper]{article}

\usepackage[utf8]{inputenc}
\usepackage[spanish]{babel}
\usepackage[T1]{fontenc}

\usepackage[
    top=3cm,
    left=4cm,
    bottom=3cm,
    right=2cm
]{geometry}

\usepackage{setspace}
\onehalfspacing

\usepackage{mathptmx} 

\usepackage{fancyhdr}
\pagestyle{fancy}
\fancyhf{}
\fancyfoot{}
\fancyhead[R]{\thepage}
\renewcommand{\headrulewidth}{0pt}

\usepackage{graphicx}
\usepackage{caption}
\captionsetup{font=small, labelfont=bf}

\usepackage[backend=biber, style=numeric, sorting=none]{biblatex}
\addbibresource{references.bib}   

\usepackage{titlesec}
\titleformat{\section}{\normalfont\bfseries\large}{\thesection.}{1em}{\MakeUppercase}
\titleformat{\subsection}{\normalfont\bfseries}{\thesubsection}{1em}{\MakeUppercase}
\titleformat{\subsubsection}{\normalfont\bfseries\normalsize}{\thesubsubsection}{1em}{\MakeUppercase}

\begin{document}
\begin{titlepage}
    \thispagestyle{empty}
    \centering
    \vspace*{2cm}
    {\LARGE \textbf{CHATBOTS CONVERSACIONALES VS. FORMULARIOS TRADICIONALES EN LA UPTC}}\\[3cm]
    {\large Edwin David Fino Velandia\\ Edward Mauricio Guatibonza Suárez\\ Fredy Alejandro Bonilla Rodríguez}\\[2cm]
    {\large Universidad Pedagógica y Tecnológica de Colombia (UPTC)}\\
    {\large Facultad de Ingeniería}\\
    {\large Escuela de Ingeniería de Sistemas y Computación}\\[2cm]
    {\large Tunja}\\
    {\large \the\year}
\end{titlepage}

\newpage

\thispagestyle{empty}
\tableofcontents
\newpage

\setcounter{page}{1}

Texto aquí...

\section{Planteamiento del Problema}

\subsection{Antecedentes}

En las instituciones de educación superior, la recolección de información ha pasado progresivamente de formatos físicos a herramientas digitales~\cite{4, 5}. A pesar de este avance, los formularios web tradicionales todavía presentan dificultades importantes, como bajos niveles de respuesta y abandono antes de completar las preguntas. Por esta razón, se ha empezado a estudiar el uso de chatbots en aplicaciones de mensajería instantánea como una alternativa que podría facilitar la participación de los estudiantes~\cite{1, 2, 3, 4, 5}.

\subsubsection*{Antecedentes de Efectividad Comparativa: Chatbots y Formularios Web}

\begin{itemize}
    \item \textbf{Xiao et al. (2020):} Realizaron un estudio de campo con cerca de 600 participantes para comparar la recolección de datos mediante un chatbot con inteligencia artificial y un formulario web tradicional~\cite{2, 3, 4, 5}. Los resultados mostraron que el chatbot logró una tasa de finalización del 54\%, mientras que el formulario web alcanzó solo el 24.2\%~\cite{2, 3, 4, 5}. Este estudio sugiere que una interacción guiada, en la que las preguntas se presentan una a una, puede motivar a las personas a completar la encuesta.

    \item \textbf{Beam (2023):} Analizó el uso de redes sociales y agentes conversacionales para reclutar participantes y recolectar información~\cite{2, 3, 4, 5}. A partir de un experimento aleatorizado aplicado a encuestas docentes, encontró que los chatbots obtuvieron mejores tasas de respuesta y datos de mayor calidad que los formularios web convencionales~\cite{2, 3, 4, 5}.

    \item \textbf{Zarouali, Araujo y Ohme (2023):} Desarrollaron una investigación longitudinal para comparar la respuesta de los participantes, la calidad de los datos y la percepción de uso entre chatbots y encuestas web~\cite{2, 3, 4, 5}. A diferencia de los estudios anteriores, no encontraron ventajas estadísticamente significativas del chatbot frente al formulario tradicional a lo largo del tiempo. Además, observaron una disminución en la valoración positiva de los usuarios en mediciones repetidas~\cite{2, 3, 4, 5}. Estos resultados muestran que la efectividad de los chatbots puede depender del contexto y que es necesario evaluarlos en escenarios específicos antes de asumir que siempre producirán mejores resultados.
\end{itemize}

Por su parte, \textbf{Beam (2023)} evaluó experimentalmente la efectividad de las redes sociales y agentes conversacionales como herramientas de reclutamiento y recolección de información~\cite{2, 3, 4, 5}. Mediante un diseño experimental aleatorizado aplicado a encuestas docentes, el estudio demostró que los chatbots alcanzaron tasas de respuesta y una calidad de datos significativamente superiores en comparación con los formularios web estáticos convencionales~\cite{2, 3, 4, 5}.

Sin embargo, \textbf{Zarouali, Araujo y Ohme (2023)} llevaron a cabo una investigación longitudinal comparando el comportamiento de respuesta, la calidad de datos y la evaluación de usuario entre chatbots y encuestas web~\cite{2, 3, 4, 5}. A diferencia de los hallazgos de Xiao et al. y Beam, esta investigación no encontró ventajas estadísticamente significativas en el uso del chatbot frente al formulario tradicional en el tiempo, e incluso registró una disminución en la valoración favorable por parte de los usuarios en mediciones repetidas~\cite{2, 3, 4, 5}. Este resultado evidencia la falta de consenso en la comunidad científica y la necesidad de evitar asumir resultados de antemano sin una contrastación local~\cite{2, 3, 4, 5}.

\subsubsection*{Antecedentes sobre Diseño de Interacción y Usabilidad}

\textbf{Kim, Lee y Gweon (2019)} \cite{kim2019} analizaron el impacto de la plataforma y el estilo conversacional en la calidad de las respuestas en encuestas en línea. Sus hallazgos señalan que los formularios web tradicionales operan como sistemas receptores pasivos al carecer de confirmación en tiempo real; esta rigidez choca con la expectativa de interactividad inmediata propia de la mensajería instantánea.

Por su parte, \textbf{Hoerger (2010)} \cite{hoerger2010} examinó la deserción de participantes en función de la longitud del cuestionario en entornos universitarios. La evidencia cuantitativa demostró que cerca del 10\% de los estudiantes abandona la prueba de inmediato al percibir una interfaz extensa, y que la tasa de deserción se incrementa aproximadamente 2 puntos porcentuales por cada 100 preguntas adicionales.

Desde la perspectiva del diseño de interfaz, \textbf{Nielsen y Molich (1990)} \cite{nielsen1990} establecieron los principios de la evaluación heurística. Sus postulados explican cómo la falta de retroalimentación inmediata, la escasa visibilidad del estado del sistema y la ausencia de indicadores de progreso en formularios estáticos generan carga cognitiva e inducen al abandono prematuro.

Asimismo, \textbf{Harley et al. (2007)} \cite{harley2007} evaluaron el uso de la mensajería móvil en la comunicación académica universitaria. Los resultados revelaron que la comunidad estudiantil otorga prioridad y revisión continua a las aplicaciones de mensajería instantánea por encima de los canales formales como el correo electrónico, donde las notificaciones tienden a postergarse e ignorarse.

En el ámbito metodológico, \textbf{Hernández-Sampieri, Fernández-Collado y Baptista-Lucio (2014)} \cite{hernandez2014} advierten que una alta tasa de deserción no aleatoria en la recolección de datos genera un sesgo de selección. Este fenómeno compromete la validez externa de los resultados e impide la representatividad de la muestra respecto a la población de estudio.

La revisión bibliográfica evidencia una discusión activa en la literatura internacional. Por un lado, diversos estudios empíricos afirman que los chatbots conversacionales elevan la tasa de finalización y la calidad de la información al integrarse al canal cotidiano del usuario \cite{xiao2020, beam2023}. Por otro lado, investigaciones longitudinales señalan que estos entornos no siempre superan a las encuestas web e incluso pueden suscitar desgaste en el participante \cite{zarouali2023}.

Esta tensión teórica fundamenta la necesidad de ejecutar un experimento comparativo en el contexto de la Universidad Pedagógica y Tecnológica de Colombia (UPTC). El estudio busca evaluar si la implementación de un chatbot conversacional en Telegram logra mitigar la brecha de canal y la fricción de usabilidad que afectan a los formularios web institucionales \cite{uptc2024}.

\subsection{Descripción de la Situación Problemática}
Texto aquí...

\subsection{Formulación del Problema}
Texto aquí...

\subsection{Consecuencias de no Investigar}
Texto aquí...

\subsection{Delimitación}
Texto aquí...

\section{Justificación}
Texto aquí...

\section{Objetivos}
\subsection{Objetivo General}
Texto aquí...

\subsection{Objetivos Específicos}
Texto aquí...

\section{Alcance de la Investigación}
Texto aquí...

\section{Formulación de Hipótesis}
Texto aquí...

\section{Variables}
Texto aquí...

\section{Marco Referencial}

\subsection{Marco Teórico Conceptual}

\subsubsection{Pruebas de Software}
Texto aquí...

\subsubsection{Automatización de Pruebas}
Texto aquí...

\subsubsection{Priorización de Pruebas}
Texto aquí...

\subsection{Marco Histórico}
Texto aquí...

\subsection{Estado Actual del Tema}
Texto aquí...

\subsection{Marco Científico y Tecnológico}

\subsubsection{Fundamentos Científicos}
Texto aquí...

\subsubsection{Tecnologías Relevantes}
Texto aquí...

\subsection{Marco Legal y Normativo}
Texto aquí... 

\section{Diseño Metodológico}
Texto aquí...

\section{Resultados y Discusión}
Texto aquí...

\section{Conclusiones y Recomendaciones}
Texto aquí...

\printbibliography[title={Bibliografía}]

\section*{Anexos}
Texto aquí...

\end{document}
